% SEGMENT OLP-0026-S01
% Part: sets-functions-relations
% Chapter: functions
% Section: partial-functions

\documentclass[../../../include/open-logic-section]{subfiles}

\begin{document}

\olfileid{sfr}{fun}{par}

\olsection{பகுதிச் சார்புகள்}

\begin{explain}
சில சமயங்களில் சார்பின் வரையறையைத் தளர்த்துவது பயனுள்ளது:
சாத்தியமான எல்லா உள்ளீடுகளுக்கும் சார்பின் வெளியீடு வரையறுக்கப்பட்டிருக்க
வேண்டும் என்ற தேவையை நீக்கலாம்.
இத்தகைய இணைப்புகள் \emph{பகுதிச் சார்புகள்} எனப்படும்.
\end{explain}

% SEGMENT OLP-0026-S02
\begin{defn}
ஒரு \emph{பகுதிச் சார்பு} $f \colon A \pto B$ என்பது
$A$ இன் ஒவ்வோர் !!{element} உடனும் $B$ இன் அதிகபட்சம்
ஓர் !!{element}[acc] இணைக்கும் ஓர் இணைப்பாகும்.
$x \in A$ உடன் $B$ இன் ஓர் உறுப்பை $f$ இணைக்குமானால்,
$f(x)$ \emph{வரையறுக்கப்பட்டுள்ளது} என்கிறோம்;
இல்லையெனில் அது \emph{வரையறுக்கப்படவில்லை} என்கிறோம்.
$f(x)$ வரையறுக்கப்பட்டிருந்தால் $f(x) \fdefined$ எனவும்,
இல்லையெனில் $f(x) \fundefined$ எனவும் எழுதுகிறோம்.
ஒரு பகுதிச் சார்பு $f$ இன் \emph{சார்பகம்} என்பது,
அது வரையறுக்கப்பட்டுள்ள $A$ இன் உட்கணமாகும்;
அதாவது $\dom{f} = \Setabs{x \in A}{f(x) \fdefined}$.
\end{defn}

% SEGMENT OLP-0026-S03
\begin{ex}
ஒவ்வொரு சார்பு $f\colon A \to B$ உம் ஒரு பகுதிச் சார்பாகவும் உள்ளது.
$A$ முழுவதிலும் வரையறுக்கப்பட்ட பகுதிச் சார்புகள்,
அதாவது இதுவரை நாம் வெறுமனே சார்புகள் என அழைத்தவை,
\emph{முழுச்} சார்புகள் என்றும் அழைக்கப்படுகின்றன.
\end{ex}

% SEGMENT OLP-0026-S04
\begin{ex}
$f(x) = 1/x$ எனக் கொடுக்கப்படும்
$f \colon \Real \pto \Real$ என்ற பகுதிச் சார்பு,
$x = 0$ என்பதில் வரையறுக்கப்படவில்லை;
மற்ற எல்லா இடங்களிலும் வரையறுக்கப்பட்டுள்ளது.
\end{ex}

% SEGMENT OLP-0026-S05
\begin{prob}
$f\colon A \pto B$ தரப்பட்டிருக்க, $g\colon B \pto
A$ என்ற பகுதிச் சார்பைப் பின்வருமாறு வரையறுப்போம்:
எந்த $y \in B$ க்கும், $f(x) = y$ என அமையும்
தனித்த $x \in A$ இருந்தால் $g(y) = x$;
இல்லையெனில் $g(y) \fundefined$.
$f$ என்பது ஒன்றுக்கொன்றானது எனில்,
அனைத்து $x \in \dom{f}$ க்கும் $g(f(x)) = x$ எனவும்,
அனைத்து $y \in \ran{f}$ க்கும் $f(g(y)) = y$ எனவும் காட்டுக.
\end{prob}

% SEGMENT OLP-0026-S06
\begin{defn}[பகுதிச் சார்பின் கோட்டுரு]
$f\colon A \pto B$ ஒரு பகுதிச் சார்பாக இருக்கட்டும்.
$f$ இன் \emph{கோட்டுரு} என்பது பின்வருமாறு வரையறுக்கப்படும்
$R_f \subseteq A \times B$ என்ற தொடர்பாகும்:
\[
R_f = \Setabs{\tuple{x,y}}{f(x) = y}.
\]
\end{defn}

% SEGMENT OLP-0026-S07
\begin{prop}
$R \subseteq A \times B$ என்பது, $Rxy$, $Rxy'$ என அமையும்
போதெல்லாம் $y = y'$ என்ற பண்பைக் கொண்டிருப்பதாகக் கொள்வோம்.
அப்போது $R$ என்பது பின்வருமாறு வரையறுக்கப்படும்
$f\colon A \pto B$ என்ற பகுதிச் சார்பின் கோட்டுருவாகும்:
$Rxy$ என அமையும் ஒரு $y$ இருந்தால் $f(x) = y$;
இல்லையெனில் $f(x) \fundefined$.
$R$ மேலும் \emph{தொடர் பண்புடையதாக}, அதாவது ஒவ்வொரு
$x \in A$ க்கும் $Rxy$ என அமையும் ஒரு $y \in B$ உள்ளது
என்ற பண்புடையதாக, இருந்தால் $f$ முழுச் சார்பாகும்.
\end{prop}

% SEGMENT OLP-0026-S08
\begin{proof}
$Rxy$ என அமையும் ஒரு $y$ இருப்பதாகக் கொள்வோம்.
$Rxy'$ என அமையும் வேறொரு $y'
\neq y$ இருந்தால், $R$ மீதான நிபந்தனை மீறப்படும்.
எனவே $Rxy$ என அமையும் ஒரு $y$ இருந்தால், அந்த $y$ தனித்ததாகும்;
ஆகவே $f$ தெளிவாக வரையறுக்கப்பட்டுள்ளது.
$R_f = R$ என்பதும், $R$ தொடர் பண்புடையது எனில் $f$ முழுச் சார்பு
என்பதும் தெளிவு.
\end{proof}

\end{document}
